\subsection{Sensitivity to the Thrust Calibration}
\label{sec:thrust_sensitivity}

One model parameter is not measured but assumed, and it deserves its
own accounting because it is the most consequential of the study. The
Riley tables supply the aerodynamics; they do not supply the map from
throttle setting to thrust. As stated in Sec.~\ref{sec:riley_aero},
$K_t$ is calibrated by requiring that full throttle sustain level
cruise at $V_{max} = 2V_s$, giving $K_t \approx 1151$\,N. Measured
against the $108$\,hp continuous rating that Riley does document
(Table~\ref{tab:aa1_params}), that calibration implies a propulsive
efficiency of $\eta \approx 0.90$ at $V_{max}$. The propeller of this
airframe is a two-blade fixed-pitch unit turning at $2600$\,rpm
(Table~\ref{tab:aa1_params}), which at $V_{max}$ places it at an advance
ratio $J = V/(nD) \approx 0.81$. Full-scale tests of two-blade
fixed-pitch propellers report, at that advance ratio, an efficiency near
$0.76$ at a representative fixed blade angle, rising to $\approx 0.85$
only along the envelope obtained by \emph{re-pitching} the blade for
each condition --- an option a fixed-pitch installation does not have
--- with an absolute peak of $\approx 0.87$ reached at
$J \approx 1.25$--$1.55$ \cite{HartmanBiermann1938}. The calibration
therefore assumes a propulsive efficiency above the peak that propeller
class attains anywhere, and roughly $18\%$ above what a fixed blade
angle delivers at the relevant operating point.
% Cifras tomadas de las cartas de NACA-TR-640 para la helice 5868-9
% (2 palas, Clark Y). Comparacion de CLASE, no de modelo: esas helices son
% de 10 ft y perfiles de 1938; la de la AA-1 es de 5.9 ft. Lo que se compara
% es el orden de magnitud del rendimiento alcanzable, no un valor puntual.

The consequence is not marginal. Flying the \emph{nominal} policy on
plants whose engine delivers $\eta \in [0.75, 0.95]$ of the modelled
thrust --- the pilot commands the same throttle fraction, the airframe
receives less --- moves the altitude loss at the canonical entry from
$6.1$ to $15.8$\,m, a factor of $2.6$. For scale, the two model
corrections that this section's aerodynamics required move the same
number by $0.4$ and $0.9$\,m. The thrust calibration dominates every
other uncertainty in the model by an order of magnitude, and no choice
of $\eta$ can be defended from the wind-tunnel data alone.

What survives the uncertainty is the comparison the paper is actually
about. Across the same range, the ordering of the recovery strategies
never changes,
\begin{equation*}
    \text{DP optimum} \;>\; \text{CAA} \;>\; \text{FAA}
    \;>\; \text{power-delayed},
\end{equation*}
and the margin of the optimum over the CAA doctrine stays between
$7.9$ and $9.2$\,m while the absolute losses swing by $9.7$\,m. The
advantage of optimal control over procedural flying is therefore a
property of the control law, not an artifact of how the propeller was
calibrated; the absolute altitude figures should be read as conditional
on $\eta \approx 0.90$, and rescaled accordingly for an airframe whose
propulsive efficiency is known.


